\begin{abstract}
We present a reinforcement-learning framework for training small language
models to reason on verifiable problems. The agent performs approximate
$Q$-learning over an LLM-induced context Markov decision process: states are
structured contexts, actions are a small set of coarse cognitive verbs the
model realizes as text, transitions are sampled from the model's responses,
and reward is anchored by a non-LLM verifier (a symbolic solver, unit-test
harness, or proof checker) rather than a learned critic. The action-value
function $\Qpsi(s, a)$ is a multi-head $Q$-network over the hidden state of a
frozen base model. Monte Carlo Tree Search enters not as the core learning
algorithm but as a policy-improvement operator: it produces locally improved
estimates $\Qmcts(s, a)$ and visit-count policies $\pimcts$ that serve as
stronger training targets for expert iteration, and it doubles as a
test-time-compute knob at inference. Keeping reward verifier-anchored rather
than learned reduces the Goodhart attack surface. We give the framework, the
architecture, and a pre-registered experimental plan with falsifiable
hypotheses (a fixed baseline grid, a FLOP-accounted compute axis, and a
Holm-Bonferroni-corrected criterion committed before training), plus a
toy-problem pilot. One round of expert iteration raises zero-shot held-out
argmax agreement by $15$ percentage points ($57\% \to 72\%$, single training
seed); a second round on naively-expanded data yields a structurally
informative regression that a controlled rebalancing test traces to
distillation-target instability rather than class dilution. We then run the
full online loop on an in-process $1.5$B substrate: a first closed-loop run
establishes the honest per-decision compute accounting, and a fully
instrumented run capturing the $Q$-loss trend and prior-sharpening signal
across iterations is currently executing, with those results pending. The
pilot is feasibility evidence; the large-scale benchmark claim is deferred to
a follow-up study.
\end{abstract}
